\chapter{Viral hemorrhagic fevers}
\index{Viral hemorrhagic fevers}

\section*{Definition and Overview}
Viral hemorrhagic fevers (VHFs) refer to a group of acute, highly infectious, and potentially life-threatening zoonotic diseases caused by several distinct families of RNA viruses. These illnesses are characterized by fever, diffuse vascular damage, multisystem organ failure, and, in severe cases, bleeding manifestations (hemorrhage). Despite the moniker ``hemorrhagic fever,'' clinical hemorrhage is not always present and is rarely the primary cause of death; rather, it is the systemic vascular dysfunction, endothelial leakage, and resulting hypovolemic or cardiogenic shock that drive patient mortality.

The clinical and public health significance of VHFs is profound. Because many of these viruses are highly transmissible in healthcare settings, carry high case-fatality rates, and have the potential for aerosol transmission in laboratory environments, they are classified as Category A bioterrorism agents and pathogens of high consequence. Outbreaks of VHFs, such as the historic 2014--2016 West African Ebola epidemic, have demonstrated their capacity to devastate weak healthcare infrastructures, cause severe socioeconomic disruption, and pose significant global health security threats. Understanding the pathophysiological mechanisms, clinical presentations, and management protocols for these diseases is critical for clinicians, public health officials, and emergency responders worldwide.

\section*{Epidemiology}
The epidemiology of VHFs is complex and closely tied to the ecology and geographic distribution of their respective animal reservoirs or arthropod vectors. These diseases are endemic in specific regions of the world where these vectors and reservoirs thrive, although international travel and global trade can occasionally lead to imported cases in non-endemic countries.

\begin{itemize}
\item \textbf{Filoviruses (Ebola and Marburg):} Primarily endemic to sub-Saharan Africa. The natural reservoir is believed to be fruit bats (family Pteropodidae). Outbreaks occur sporadically when humans come into contact with infected bats or through the consumption and handling of ``bushmeat'' (non-human primates, forest antelopes). Case fatality rates for Ebola virus disease (EVD) can range from 25\% to 90\%, depending on the viral species (e.g., Zaire ebolavirus is highly lethal) and the quality of supportive care.
\item \textbf{Arenaviruses (Lassa Fever and South American Hemorrhagic Fevers):} Lassa fever is highly endemic in West Africa, including Nigeria, Sierra Leone, Liberia, and Guinea, causing an estimated 100,000 to 300,000 cases annually with approximately 5,000 deaths. The reservoir is the multimammate rat (\textit{Mastomys natalensis}). The South American hemorrhagic fevers (Junin, Machupo, Guanarito, Sabia, and Lujo) are caused by rodent-borne arenaviruses endemic to specific agricultural regions of Argentina, Bolivia, Venezuela, and Brazil, respectively.
\item \textbf{Flaviviruses (Dengue and Yellow Fever):} Yellow fever is endemic in tropical areas of Africa and Central and South America. While many infections are mild, severe cases can lead to classic hemorrhagic manifestations and liver failure. Dengue is the most widespread vector-borne viral disease globally, transmitted by \textit{Aedes} mosquitoes, with billions of people at risk in tropical and subtropical regions. Severe dengue (dengue hemorrhagic fever/dengue shock syndrome) is a major cause of hospitalization and death among children in Asia and Latin America.
\item \textbf{Bunyaviruses (Crimean-Congo Hemorrhagic Fever, Rift Valley Fever, and Hantaviruses):} Crimean-Congo hemorrhagic fever (CCHF) is the most geographically widespread tick-borne VHF, endemic in Eastern Europe, the Mediterranean, northwestern China, Central Asia, and parts of Africa and the Middle East, transmitted by \textit{Hyalomma} ticks. Rift Valley fever (RVF) is primarily zoonotic, affecting livestock and humans in Africa and the Arabian Peninsula, transmitted by mosquitoes or contact with infected animal blood. Hantaviruses, causing Hemorrhagic Fever with Renal Syndrome (HFRS), are distributed across Eurasia and are transmitted through contact with rodent excreta.
\end{itemize}

Affected populations typically include rural communities, agricultural workers, foresters, and individuals living in substandard housing with poor rodent control. Healthcare workers represent an exceptionally high-risk group during outbreaks due to direct, unprotected contact with infectious bodily fluids during patient care.

\section*{Aetiology and Pathophysiology}
The etiologic agents of VHFs are enveloped RNA viruses belonging to four distinct taxonomic families: Filoviridae, Arenaviridae, Flaviviridae, and Bunyavirales (which includes Phenuiviridae, Hantaviridae, and Nairoviridae).

The underlying pathophysiology of VHFs centers on viral tropism for endothelial cells, macrophages, and dendritic cells, leading to a cascade of immunopathological events:

\begin{itemize}
\item \textbf{Filoviridae (Ebola and Marburg):} These filoviruses infect monocytes, macrophages, and dendritic cells, triggering the systemic release of pro-inflammatory cytokines, chemokines, and vasoactive mediators (the ``cytokine storm''). This cytokine cascade leads to widespread endothelial cell activation and dysfunction, increasing vascular permeability and inducing systemic vasodilation. Concurrently, infected cells express tissue factor, initiating the extrinsic coagulation cascade and culminating in disseminated intravascular coagulation (DIC), which depletes clotting factors and platelets, leading to microvascular thrombosis and hemorrhagic diathesis.
\item \textbf{Arenaviridae (Lassa and South American VHFs):} These viruses enter host cells via specific cellular receptors, such as alpha-dystroglycan. Arenaviruses actively suppress the host's innate immune response by inhibiting type I interferon production, allowing unrestricted viral replication. The clinical manifestations result from direct cellular injury and systemic capillary leakage caused by inflammatory cytokines, rather than direct viral lysis of endothelial cells.
\item \textbf{Flaviviridae (Dengue and Yellow Fever):} In severe dengue, the phenomenon of antibody-dependent enhancement (ADE) plays a critical role. Non-neutralizing antibodies from a prior infection bind to a different serotype of dengue virus during a secondary infection, facilitating viral entry into host macrophages via Fc receptors. This leads to massive activation of T-cells and the release of vasoactive cytokines, resulting in transient but severe plasma leakage. In Yellow Fever, the virus displays strong hepatotropism, causing extensive lytic necrosis of hepatocytes (Councilman bodies), leading to severe jaundice, coagulopathy due to impaired synthesis of clotting factors, and renal failure.
\item \textbf{Bunyavirales (CCHF, RVF, Hantaviruses):} CCHF virus directly targets hepatocytes and endothelial cells. The virus impairs the host's immune response through the suppression of interferon pathways, while also triggering TNF-alpha and other cytokines that damage the endothelial barrier. Hantaviruses primarily target the endothelial cells of the renal and pulmonary vasculature, causing increased capillary permeability and fluid shifts without causing significant cytopathic destruction of the vessels.
\end{itemize}

Ultimately, all VHFs share a common pathway of endothelial damage, capillary leakage, intravascular volume depletion, profound hypotension, disseminated intravascular coagulation, and multi-organ dysfunction syndrome (MODS).

\section*{Clinical Features}
The clinical presentation of VHFs is highly variable, ranging from mild, self-limiting febrile illnesses to rapidly progressive, fatal shock. The incubation period varies by virus, typically ranging from 2 to 21 days. The clinical course can generally be divided into distinct stages:

\begin{itemize}
\item \textbf{Prodromal (Early) Phase:} This phase is marked by the sudden onset of non-specific systemic symptoms that mimic common tropical illnesses. Symptoms include:
  \begin{itemize}
  \item High fever and chills
  \item Severe headache, retro-orbital pain, and photophobia
  \item Generalized myalgia, arthralgia, and lower back pain
  \item Extreme fatigue, malaise, and physical weakness
  \item Gastrointestinal disturbances, including anorexia, nausea, vomiting, abdominal pain, and watery diarrhea
  \item Sore throat, dry cough, and conjunctival injection (red eyes)
  \end{itemize}
\item \textbf{Hemorrhagic and Vascular Leak Phase:} Occurs in severe cases, typically 5 to 7 days after symptom onset. Signs of increased vascular permeability and coagulopathy become prominent:
  \begin{itemize}
  \item Mucosal bleeding, including epistaxis (nosebleeds), gingival bleeding (bleeding gums), and subconjunctival hemorrhage
  \item Gastrointestinal hemorrhage presenting as hematemesis (vomiting blood) or melena (black, tarry stools)
  \item Cutaneous manifestations, such as petechiae, purpura, ecchymoses, and bleeding from venipuncture sites
  \item Hematuria and metrorrhagia
  \item Postural hypotension, tachycardia, and signs of plasma leakage (ascites, pleural effusion, and generalized edema)
  \item Maculopapular rash, particularly prominent in filovirus infections, appearing on the trunk and spreading to the extremities
  \end{itemize}
\item \textbf{Late/Organ Failure Phase:} In fatal cases, the illness progresses rapidly to:
  \begin{itemize}
  \item Profound hypovolemic and distributive shock
  \item Acute kidney injury (AKI) presenting as oliguria or anuria
  \item Hepatic encephalopathy, jaundice, and severe metabolic acidosis
  \item Neurological involvement, including confusion, agitation, tremors, seizures, obtundation, and coma
  \item Acute respiratory distress syndrome (ARDS) or tachypnea
  \item Cardiovascular collapse and refractory hypotension
  \end{itemize}
\end{itemize}

During recovery, survivors may experience a prolonged convalescence phase characterized by severe weakness, arthralgias, hair loss, desquamation of the skin, uveitis, and various neuropsychiatric symptoms.

\section*{Differential Diagnosis}
Given the non-specific nature of the early symptoms, the differential diagnosis for VHFs is broad and encompasses numerous infectious and non-infectious conditions. Accurate differentiation is critical for initiating appropriate isolation measures and supportive therapies.

\begin{itemize}
\item \textbf{Malaria:} Specifically \textit{Plasmodium falciparum} malaria, which is co-endemic in many VHF regions. It presents with fever, splenomegaly, anemia, and thrombocytopenia, and can lead to cerebral malaria and multi-organ failure. It is ruled out via thick and thin peripheral blood smears or rapid diagnostic tests (RDTs).
\item \textbf{Typhoid Fever:} Caused by \textit{Salmonella enterica} serovar Typhi. Characterized by prolonged fever, abdominal pain, bradycardia (Faget's sign), rose spots, and splenomegaly. Confirmed via blood, stool, or bone marrow cultures.
\item \textbf{Leptospirosis:} A zoonotic bacterial infection transmitted via water contaminated with rodent urine. Severe leptospirosis (Weil's disease) features jaundice, renal failure, and hemorrhage, closely mimicking Yellow Fever or CCHF. Diagnosed by serology, PCR, or microagglutination testing.
\item \textbf{Sepsis and Meningococcemia:} Bacterial septic shock, particularly due to \textit{Neisseria meningitidis}, can cause rapid-onset purpura fulminans, petechial rash, DIC, and circulatory collapse. Blood cultures and cerebrospinal fluid analysis are key diagnostic tools.
\item \textbf{Rickettsial Diseases:} Such as Rocky Mountain spotted fever or epidemic typhus. These present with fever, headache, myalgia, and a macular or petechial rash. Confirmed through serology or PCR.
\item \textbf{Influenza and Severe Acute Respiratory Infections:} During the prodromal phase, respiratory viruses can mimic the initial presentation of VHFs.
\item \textbf{Hematologic Malignancies and Autoimmune Disorders:} Acute leukemia or severe systemic lupus erythematosus can present with cytopenias, fever, and bleeding manifestations.
\end{itemize}

\section*{When to Seek Medical Care}
Suspicion of a viral hemorrhagic fever constitutes a medical emergency. Immediate medical attention must be sought if an individual develops a sudden, unexplained high fever accompanied by severe headache, muscle aches, vomiting, diarrhea, or any unusual bruising or bleeding, particularly if they have:
\begin{itemize}
\item Traveled to a VHF-endemic region (e.g., sub-Saharan Africa, parts of South America, rural Asia, or the Middle East) within the past 21 days.
\item Had known or suspected contact with bats, rodents, or livestock in endemic areas.
\item Participated in outdoor activities in areas where tick or mosquito vectors are highly active.
\item Provided care for an individual who is ill with unexplained febrile or hemorrhagic symptoms.
\end{itemize}

If a patient exhibits life-threatening symptoms, emergency medical services must be contacted immediately. Do not visit a clinic or hospital without notifying them in advance of your travel history and symptoms to prevent potential transmission. Emergency call numbers include:
\begin{itemize}
\item \textbf{local emergency services} for the affected region
\item \textbf{911} for the United States and Canada
\item \textbf{999} for the United Kingdom
\item \textbf{112} for the European Union and internationally from mobile networks
\end{itemize}
When calling, explicitly state the suspicion of a highly infectious viral pathogen so that responding paramedic and healthcare teams can wear appropriate personal protective equipment (PPE) before contact.

\section*{Diagnosis and Investigations}
Diagnosing VHFs requires a high index of clinical suspicion combined with a detailed travel and exposure history. Because handling specimens from suspected VHF patients poses severe biosafety risks, laboratory testing must be conducted under strict Biosafety Level 4 (BSL-4) containment or using specific validated point-of-care assays with strict biosafety precautions.

\begin{itemize}
\item \textbf{Specific Diagnostic Assays:}
  \begin{itemize}
  \item \textbf{Reverse Transcription-Polymerase Chain Reaction (RT-PCR):} The gold standard for early diagnosis. RT-PCR detects viral RNA in blood or tissues during the acute phase of the illness, typically within the first 3 to 10 days of symptom onset.
  \item \textbf{Antigen-Capture Enzyme-Linked Immunosorbent Assay (ELISA):} Detects viral antigens (such as nucleoproteins) in the blood. Highly useful in the early stages of disease when viral load is high.
  \item \textbf{Serology (IgM and IgG ELISA):} Used to detect antibodies. IgM antibodies appear during the late acute phase (around day 7--10) and indicate recent infection. IgG antibodies appear later and persist for years, useful for epidemiological surveillance and convalescent confirmation.
  \item \textbf{Viral Isolation:} Culturing the virus in cell lines. This is restricted exclusively to BSL-4 facilities due to the extreme risk of laboratory-acquired infection.
  \end{itemize}
\item \textbf{General Supportive Investigations:}
  \begin{itemize}
  \item \textbf{Complete Blood Count (CBC):} Characterized by marked leukopenia (often with a left shift) and severe thrombocytopenia. Hematocrit may be elevated initially due to hemoconcentration from plasma leakage.
  \item \textbf{Coagulation Profile:} Prolonged prothrombin time (PT), activated partial thromboplastin time (aPTT), elevated D-dimers, and decreased fibrinogen levels, indicative of DIC.
  \item \textbf{Comprehensive Metabolic Panel (CMP):} Reveals elevated liver transaminases (aspartate aminotransferase [AST] is typically significantly higher than alanine aminotransferase [ALT] in filovirus and Lassa infections), elevated blood urea nitrogen (BUN) and creatinine (indicating renal impairment), and electrolyte disturbances (hypokalemia, hyponatremia).
  \item \textbf{Urinalysis:} Often shows proteinuria and microscopic hematuria.
  \end{itemize}
\end{itemize}

\section*{Management}
The primary management of VHFs relies on aggressive, high-quality supportive care, as specific antiviral therapies are limited. Management must be carried out in specialized isolation units with strict infection prevention and control (IPC) measures to protect healthcare staff and the community.

\begin{itemize}
\item \textbf{Fluid and Electrolyte Management:} The cornerstone of treatment is maintaining intravascular volume and correcting electrolyte imbalances. This involves:
  \begin{itemize}
  \item Oral rehydration solutions (ORS) for patients who can tolerate oral intake.
  \item Tailored intravenous fluids (crystalloids or colloids) to counter plasma leakage and hypovolemia, guided by hemodynamic monitoring (heart rate, blood pressure, urine output) to avoid volume overload and pulmonary edema.
  \item Maintenance of electrolyte balance, particularly monitoring and replacing potassium, magnesium, and calcium.
  \item Vasopressors (e.g., norepinephrine) in patients with refractory shock despite adequate fluid resuscitation.
  \end{itemize}
\item \textbf{Pharmacological Interventions:}
  \begin{itemize}
  \item \textbf{Monoclonal Antibodies:} Specifically developed and approved for Zaire ebolavirus. Single-dose intravenous infusions of Ebanga (ansuvimab-zykl) or Inmazeb (atoltivimab, maftivimab, and odesivimab-ebgn) have been shown to significantly reduce mortality when administered early in the course of the disease.
  \item \textbf{Antiviral Agents:} Ribavirin is an antiviral drug that has shown efficacy against certain Arenaviruses (Lassa fever) and Bunyaviruses (CCHF) if administered early in the course of infection. It is not effective against filoviruses or flaviviruses.
  \item \textbf{Analgesics and Antipyretics:} Acetaminophen (paracetamol) is used for fever and pain control. Non-steroidal anti-inflammatory drugs (NSAIDs) such as ibuprofen, aspirin, and naproxen are strictly contraindicated due to their antiplatelet effects, which exacerbate bleeding risks.
  \item \textbf{Empiric Antimicrobial Therapy:} Broad-spectrum antibiotics (e.g., ceftriaxone) and antimalarials (e.g., artemether-lumefantrine) are often administered concurrently while waiting for definitive laboratory results, as bacterial sepsis and malaria are common co-infections.
  \end{itemize}
\item \textbf{Coagulation and Hemostasis Support:} In patients with clinical hemorrhage and coagulopathy, transfusion of packed red blood cells, fresh frozen plasma, cryoprecipitate, or platelets may be indicated, although this must be balanced against the risk of fluid overload.
\item \textbf{Renal Replacement Therapy:} Peritoneal dialysis or continuous renal replacement therapy (CRRT) may be utilized in specialized centers to manage acute kidney injury and fluid overload, provided strict biosafety protocols are maintained.
\end{itemize}

\section*{Complications}
VHFs are associated with severe, systemic complications that carry high morbidity and mortality:

\begin{itemize}
\item \textbf{Hypovolemic and Distributive Shock:} Caused by massive plasma leakage, gastrointestinal fluid losses, and systemic vasodilation, leading to inadequate tissue perfusion.
\item \textbf{Disseminated Intravascular Coagulation (DIC):} Widespread microvascular thrombosis leading to organ ischemia, coupled with consumption of clotting factors and platelets, leading to catastrophic hemorrhage.
\item \textbf{Multi-Organ Dysfunction Syndrome (MODS):} Concurrent failure of the renal, hepatic, cardiovascular, and respiratory systems.
\item \textbf{Acute Kidney Injury (AKI):} Resulting from renal hypoperfusion (prerenal azotemia), direct viral cytopathology, or acute tubular necrosis.
\item \textbf{Encephalopathy and Neurological Deficits:} Seizures, cerebral edema, coma, and focal neurological deficits. In survivors, long-term complications can include sensorineural hearing loss (common in Lassa fever survivors), uveitis, peripheral neuropathy, and cognitive dysfunction.
\item \textbf{Bacterial Superinfections:} Secondary bacterial pneumonia or bacteremia due to damaged mucosal barriers and immune suppression.
\item \textbf{Post-VHF Syndrome:} A chronic, debilitating syndrome in survivors characterized by chronic fatigue, severe myalgias, arthralgias, headaches, memory loss, and ocular manifestations (e.g., chorioretinitis).
\end{itemize}

\section*{Prognosis and Outcomes}
The prognosis for individuals diagnosed with VHFs depends heavily on the specific viral pathogen, the viral load at presentation, the patient's age and baseline health status, and the timing and quality of supportive care.

For Ebola virus disease (Zaire ebolavirus), the case-fatality rate historical average is approximately 50\%, but it can reach 90\% without supportive care. Early administration of monoclonal antibodies and aggressive fluid resuscitation can reduce the fatality rate to under 10--20\%. Marburg virus disease carries a similarly high mortality rate of 24--88\%.

Lassa fever has an overall case-fatality rate of approximately 1\% to 2\%, but among patients hospitalized with severe disease, the mortality rate rises to 15--20\%. Pregnant women, particularly in the third trimester, have an exceptionally poor prognosis, with maternal and fetal mortality rates exceeding 80\%.

Severe Dengue (DHF/DSS) has a mortality rate of less than 1\% when managed with appropriate fluid resuscitation, but this can exceed 20\% if untreated. Yellow Fever carries a case-fatality rate of 20--50\% in patients who enter the toxic phase of hepatic and renal failure.

Survivors often require months to fully recover. The viral genome can persist in immunologically privileged sites (such as the semen, eyes, and central nervous system) for months or even years, posing a risk of sexual transmission or late relapse (e.g., Ebola relapse presenting as meningitis).

\section*{Prevention and Self-Care}
Preventing VHFs requires a multi-faceted approach involving vector control, public health education, vaccine distribution, and strict infection control practices.

\begin{itemize}
\item \textbf{Vaccination:}
  \begin{itemize}
  \item \textbf{Ebola Vaccine:} Ervebo (rVSV-ZEBOV) is a highly effective, live-attenuated recombinant vaccine approved for the prevention of Zaire ebolavirus. It is used in ring-vaccination strategies during outbreaks to protect contacts of confirmed cases and frontline workers.
  \item \textbf{Yellow Fever Vaccine:} A highly effective, single-dose live-attenuated vaccine (17D) that provides lifelong immunity for over 99\% of vaccinees. It is mandatory for travel to and from endemic regions under the International Health Regulations.
  \item \textbf{Dengue Vaccine:} Vaccines such as Qdenga and Dengvaxia are approved in various countries, with specific recommendations based on serostatus and local epidemiology.
  \item \textbf{Argentine Hemorrhagic Fever Vaccine:} Candid \#1 vaccine is highly effective and widely used in endemic agricultural regions of Argentina.
  \end{itemize}
\item \textbf{Vector and Reservoir Control:}
  \begin{itemize}
  \item Avoid contact with rodents, particularly in West Africa (Lassa fever). Keep food stored in rodent-proof containers, discard garbage away from homes, and avoid handling or consuming wild rodents.
  \item Prevent mosquito and tick bites by using insect repellent containing DEET, picaridin, or IR3535; wearing long-sleeved shirts and long pants; utilizing permethrin-treated clothing and bed nets; and eliminating standing water around residences to reduce mosquito breeding.
  \item Avoid contact with fruit bats and caves where they roost.
  \item Wear gloves and protective clothing when handling livestock or veterinary tissues to prevent transmission of CCHF and RVF.
  \end{itemize}
\item \textbf{Infection Prevention and Control in Healthcare Settings:}
  \begin{itemize}
  \item Strict adherence to standard, contact, and droplet precautions.
  \item Use of appropriate personal protective equipment (PPE), including double gloves, fluid-resistant gowns, face shields, and respirators (N95 or higher).
  \item Immediate isolation of suspected cases in single rooms with dedicated sanitary facilities.
  \item Safe handling, transport, and disposal of laboratory specimens and infectious waste.
  \item Safe and dignified burial protocols for deceased patients, as bodies remain highly infectious.
  \end{itemize}
\item \textbf{Self-Monitoring and Travel Safety:} Travelers to endemic regions should monitor their health for 21 days post-return. If a fever or other symptoms develop, they must isolate immediately and seek medical attention, disclosing their travel history.
\end{itemize}

\section*{Sources and Further Reading}
\begin{itemize}
\item World Health Organization (WHO). \textit{Viral Haemorrhagic Fevers}. Available at: \url{https://www.who.int/health-topics/viral-haemorrhagic-fevers}
\item Centers for Disease Control and Prevention (CDC). \textit{Viral Hemorrhagic Fevers (VHFs)}. Available at: \url{https://www.cdc.gov/vhf/}
\item European Centre for Disease Prevention and Control (ECDC). \textit{Factsheet about Viral Haemorrhagic Fevers}. Available at: \url{https://www.ecdc.europa.eu/en/viral-haemorrhagic-fevers}
\item Gavi, the Vaccine Alliance. \textit{Preventing epidemics of viral haemorrhagic fevers}. Available at: \url{https://www.gavi.org/}
\item MedlinePlus, National Library of Medicine. \textit{Hemorrhagic Fevers}. Available at: \url{https://medlineplus.gov/hemorrhagicfevers.html}
\item The Lancet Infectious Diseases. \textit{Global burden and management of viral hemorrhagic fevers}. Available at: \url{https://www.thelancet.com/journals/laninf/home}
\end{itemize}